% ---------------------------------------------------------------------------
% Cost-aware wall-clock study (Sept 2026 revision). Everything in this file is
% new material and is rendered in red through the enclosing colour group.
% Numbers are pulled from costaware_tables.tex (generated by make_tables.py).
% ---------------------------------------------------------------------------
\section{Cost-Aware Routing and Wall-Clock Evaluation} \label{sec:wallclock}

The experiments of \Cref{sec:experiments} compare methods per iteration, which favors methods whose iterations are expensive. This appendix evaluates the end-to-end \emph{wall-clock time to a target accuracy} of the learned router against HINTS, the classical solver alone and the classical methods a practitioner would reach for on these problems (geometric multigrid and preconditioned Krylov methods), for all five classical solvers and three PDEs, on $128^2$, $256^2$ and $512^2$ grids, using the cost-aware instantiation of the greedy rule introduced at the end of \Cref{sec:greedy}. Besides the Poisson and convection--diffusion equations of \Cref{sec:experiments} we add \emph{anisotropic diffusion}, $-\epsilon\,\partial_{x_1}^2 u - \partial_{x_2}^2 u = f$ with $\epsilon = 10^{-2}$ (diffusion tensor $\mathrm{diag}(\epsilon, 1)$), on the same periodic domain with the same forcing distribution. It is the standard model problem on which point-relaxation methods fail to smooth: an error component oscillatory in $x_1$ and smooth in $x_2$ is damped by only $(1-\epsilon)/(1+\epsilon) \approx 0.98$ per Jacobi sweep, and geometric multigrid with point smoothing loses its grid-independent convergence (line relaxation or semi-coarsening would be needed). \Cref{sec:wallclock_protocol} describes the protocol; \Cref{sec:wc_poisson,sec:wc_conv,sec:wc_aniso} collect all results per equation (every grid, every tolerance, every baseline, with significance tests and seed trials); \Cref{sec:wallclock_ens} studies ensembles of one, two and three classical solvers; \Cref{sec:wallclock_overheads} accounts for overheads and training cost; and \Cref{sec:wallclock_assumptions} verifies numerically that the assumptions of \Cref{th:suboptgreedy}, Propositions~\ref{th:weaklyalphasupermodular}--\ref{th:simultaneoussupermodular} and \Cref{th:consistency} hold in these settings.

\subsection{Protocol} \label{sec:wallclock_protocol}

\paragraph{Fast, matched implementations.} A wall-clock comparison is meaningful only if the classical baseline is not handicapped by implementation. All solvers are therefore re-implemented for the periodic constant-coefficient setting as $O(N^2)$ operations: Jacobi and damped Jacobi as 5-point stencil sweeps, Gauss--Seidel, SOR and symmetric GS via cached sparse triangular factorizations, geometric multigrid as V(2,2)-cycles with lexicographic Gauss--Seidel smoothing, full-weighting restriction, bilinear prolongation, rediscretised coarse operators and a direct solve on the $32\times 32$ coarsest grid (error contraction $\approx 0.03$ per cycle on the isotropic equations). Their one-step iterates match the dense implementations used in the rest of the paper to $10^{-12}$ in double precision. Reference solutions are computed by FFT diagonalization of the same discrete operator, i.e., they are exact solutions of the discrete system, so true errors are available for evaluation at no charged cost; the FFT solve is used only as this reference (it is a direct method that exists because the test problems have constant coefficients and periodic boundaries, and is not a competitor of the iterative methods compared here). All timings are single-threaded CPU measurements (Apple M4 Pro, numpy/scipy/PyTorch, one thread) on an otherwise idle machine.

\paragraph{Corrector.} The neural corrector is a DeepONet applied to the residual $r = f_h - \mathcal{L}_h^a u^{(t)}_h$ and evaluated on a $64\times 64$ sensor grid, i.e., on the coarser grid obtained by band-limited (FFT) restriction to the modes $|k_1|,|k_2| \le 31$ at every grid size (for the anisotropic equation the sensor grid is $N\times N/4$, i.e., full resolution along $x_1$, so that the band contains every mode the point smoothers cannot damp; the sensor grid is the only PDE-specific hyperparameter of the pipeline); its output is prolongated back to the fine grid by exact trigonometric interpolation, so the correction is band-limited by construction and never injects error into the high-frequency band that the classical smoother owns (\Cref{fig:ca_predictions}). The trunk net is the fixed orthonormal Fourier basis of that band ($p = 3969$ basis functions, as in POD-DeepONet \citep{lu2022comprehensive}), and the branch net is a linear layer on the $4096$ sensor values; it is fit by least squares on $32{,}000$--$64{,}000$ exact (residual, error) pairs of the discrete operator drawn from a mixture of spectral shapes (forcing-like, residual-like, white and randomly tilted spectra), and reaches a validation relative error of $\approx 2\times 10^{-6}$ on every mode of the band. The corrector is applied scale-equivariantly, $C_{\mathrm{NO}}(r) = \|Rr\|\,\mathcal{G}_\theta(Rr/\|Rr\|)$, which preserves $C_{\mathrm{NO}}(0) = 0$ (Proposition~\ref{th:weaklyalphasupermodular}). Its cost is dominated by a $4096\times 4096$ matrix--vector product and is therefore almost independent of $N$ ($0.65$\,ms at $128^2$, $3.6$\,ms at $512^2$, \Cref{tab:ca_overheads}), i.e., about $12$ Jacobi or $3$ GS sweeps at $128^2$ and about $2.5$ Jacobi sweeps at $512^2$ (\Cref{tab:ca_costs_poisson,tab:ca_costs_conv,tab:ca_costs_aniso}). The same corrector is used by HINTS and by all routing policies.

\paragraph{Cost-equalised macro-actions.} Per-iteration costs $c_j$ (residual evaluation plus the update) are measured once per pairing and grid on the benchmark machine (minimum over seven repeated blocks of the block median, cached), the unit of cost $u$ is one corrector call, and every operation is applied $m_j = \max\{1,\mathrm{round}(u/c_j)\}$ times per decision; for the five stationary solvers this equalises costs to within $8\%$. An operation dearer than the unit (a multigrid cycle) is applied once and compared per unit of cost, $\|e_j\|^{u/(m_j c_j)}$; the surrogate weights $\tilde c_j$ use the same exponent. Ensembles reuse the cached pairwise costs of their members, so every member has the same macro-action size in the ensemble as in its pairwise run. The cost-aware oracle runs \Cref{alg:greedy} on these macro-actions with the true error; the \emph{greedy oracle} rows run the cost-agnostic \Cref{alg:greedy} on single operations, as in the main text.

\paragraph{Lightweight router.} The learned router selects macro-actions from $6 + K$ scalar features, all $O(1)$ given the residual norm that any stopping test computes: the log relative residual, its per-iteration change over the last macro-action and an exponential moving average thereof, the log epoch index, a one-hot encoding of the previous macro-action, and the log counts of corrector calls so far and of iterations since the last call. It is a two-hidden-layer MLP (64 units) evaluated in numpy, costing $\approx 5\,\mu$s per decision independently of $N$. It is trained with the surrogate of \Cref{section:surrogate}, with the per-state weights $\tilde c_k$ normalised by $\sum_k \tilde c_k$ (which does not change the Bayes decision), on states visited by oracle rollouts with $\varepsilon$-exploration (probability $0.1$), followed by two to three DAgger rounds \citep{ross2011reduction} in which the current router's own rollouts are relabelled by the oracle---the batch analogue of the scheduled sampling used in \Cref{sec:training details}. Training uses $32$--$128$ instances per round (fewer on the larger grids) drawn from the same GRF distribution as the test set but with a different seed; the router of every pairing is trained in seconds to a few minutes. Ensemble routers, which must additionally learn to choose among classical members, use a larger imitation budget ($256$ rollouts at $128^2$ and $128$ at $256^2$, six DAgger rounds, hidden width $128$, $300$ epochs): with the pairwise recipe the router over $\{\text{Jacobi}, \text{Jacobi (0.67)}\}$ imitated its oracle imperfectly at tight tolerances (convection--diffusion, $128^2$, $\varepsilon = 10^{-8}$: $3.34$ against the oracle's $2.67$ work-unit ms), which the larger budget removes ($2.70$). At $512^2$ the routers are trained from $32$ oracle rollouts of at most $400$ iterations (stopping error $10^{-8}$); for the convection--diffusion Jacobi pairing, where the first such router missed $h^2$ within the cap on $4$ of $16$ instances, the reported router uses $64$ rollouts of at most $800$ iterations---the only per-cell change to the recipe.

\paragraph{Baselines and metric.} HINTS applies the corrector every $\tau$-th iteration; we report $\tau = 25$ (the schedule used in the main text and in \citet{zhang2024blending}) and the best of $\tau \in \{5, 10, 25, 50\}$ selected \emph{per cell and per tolerance} on the test instances themselves (an optimistic baseline). The classical baselines without corrector run on the same stencil with the same accounting: multigrid V(2,2) alone; conjugate gradients, CG preconditioned by symmetric GS and by a symmetric V-cycle for Poisson; BiCGSTAB, multigrid-preconditioned BiCGSTAB and restarted GMRES(20) for convection--diffusion (GMRES is accounted at the granularity of its restart cycles and was not run at $512^2$, where it was two orders of magnitude slower than the router on the smaller grids). They are not run for anisotropic diffusion, where point-smoothed multigrid loses its grid-independent contraction ($\approx 0.85$ per cycle) and would not be informative. For each test instance ($64$ at $128^2$, $32$ at $256^2$, $16$ at $512^2$) we record, after every iteration, the true relative error $\|e^{(t)}_h\|/\|u_h\|$ (untimed, benchmark-only) and the cumulative charged wall-clock time of a replay that executes only the chosen operations (residual, update, and---for the learned router---every feature and decision computation; oracle decisions are not charged). We report the median first time at which the error drops below a tolerance $\varepsilon$, paired per-instance median speedups, and one-sided Wilcoxon signed-rank tests on the paired log time ratios (alternative: the router is faster); a run that does not reach $\varepsilon$ within the iteration cap ($60{,}000$ at $128^2$, $40{,}000$ at $256^2$, $8{,}000$ at $512^2$) is marked $\dagger$, reported as a lower bound when it is the majority, counted against the router when it is the router's, and entered at the time it spent up to the cap on the baseline side of the tests (a lower bound, hence conservative). We highlight $\varepsilon = h^2$: the discretization error of the second-order scheme scales as $O(h^2)$, so driving the algebraic error far below $h^2$ does not improve the computed solution; we also report $10^{-3}$ and $10^{-8}$, and the iteration-based metrics of the main text (final relative error and AUC $=\sum_{t=1}^{T}\|e^{(t)}_h\|/\|u_h\|$ over $T = \caT$ iterations) on the $128^2$ grid. Every router was additionally retrained from scratch with five independent seeds (different training instances, exploration and initialisation; the corrector and the test set are fixed) and re-run on the test set; because those trials ran in separate sessions, whose live timings are not comparable with the main run's, they are compared in \emph{work units} (measured per-iteration costs times executed operations, plus $5\,\mu$s per decision), which are determined by the decision sequence alone.

\begin{figure}[H]
\color{red}
    \centering
    \newfig[0.93\linewidth]{neurips_images/ca_deeponet_predictions.png}
    \caption{One-shot corrector predictions on the $128\times128$ grid (two test instances per equation): forcing, exact discrete solution, the DeepONet correction applied to the initial residual ($u^{(0)}_h = 0$), and the remaining error, which is confined to the modes outside the corrector's band (the relative error of the prediction is $10^{-6}$--$10^{-5}$ on the isotropic equations and $10^{-6}$ on anisotropic diffusion, whose solution is dominated by modes that are smooth in $x_2$).}
    \label{fig:ca_predictions}
\end{figure}

\begin{figure}[H]
\color{red}
    \centering
    \newfig[0.98\linewidth]{neurips_images/ca_convergence.png}
    \caption{Relative error versus charged wall-clock time on the $128\times128$ grid for one representative test instance of every pairing (columns) and equation (rows); the dashed line is $\varepsilon = h^2$. The learned router and the cost-aware oracle are indistinguishable.}
    \label{fig:ca_convergence}
\end{figure}

\subsection{Poisson} \label{sec:wc_poisson}

\Cref{tab:ca_time_poisson} collects the median time to $10^{-3}$, $h^2$ and $10^{-8}$ for every pairing, every policy and the classical baselines on all three grids; \Cref{tab:ca_speed_poisson} the paired speedups of the learned router with their $p$-values; \Cref{fig:ca_usage_poisson,tab:ca_usage_poisson} where the iterations go; \Cref{tab:ca_costs_poisson} the measured costs and macro-action sizes; \Cref{tab:ca_auc_poisson} the iteration-based metrics; and \Cref{tab:ca_seeds_poisson} the seed trials.

\paragraph{Uniform gains and their mechanism.} On the $128^2$ grid the learned router is the fastest deployable method at $\varepsilon = h^2$ in every pairing: hundreds to thousands of times faster than the solver alone, several times faster than HINTS ($\tau{=}25$) and still faster than the best fixed period chosen per cell and per tolerance on the test set, with all paired $p < 0.01$ (\Cref{tab:ca_speed_poisson}); at $10^{-8}$ the margins over the best fixed period shrink but remain significant in every pairing, and the best period itself varies across pairings and tolerances ($\tau{=}5$ at truncation level, $\tau \in \{5, 10, 50\}$ at $10^{-8}$), so no single fixed schedule is competitive everywhere. \Cref{tab:ca_usage_poisson} explains the numbers: with a corrector that removes the smooth band to $10^{-6}$ in one call, the cost-aware oracle and the router call it first and reach $h^2$ after one or two further sweeps, i.e., in about two iterations, whereas HINTS reaches the same accuracy only at its first scheduled call (iteration $\tau$), having spent $\tau{-}1$ sweeps on an error that is almost entirely smooth. Below $h^2$ the remaining error is the high-frequency band that the corrector cannot see; the router therefore stops calling it (one to three calls in total down to $10^{-8}$; \Cref{fig:ca_usage_poisson}), whereas a fixed schedule keeps paying for a call every $\tau$ iterations. The cost-agnostic greedy oracle (\Cref{alg:greedy} on single operations) makes the same decisions as its cost-aware counterpart at $128^2$, because after the first call the corrector is never better than a single sweep on the remaining high-frequency error; the learned router reproduces the cost-aware oracle's decisions almost exactly (same median iteration and call counts, median times within a few percent, the difference being its own decision cost and timer noise), and the five retrained routers reproduce the main router's decision sequence to $h^2$ on essentially every test instance (\Cref{tab:ca_seeds_poisson}).

\paragraph{Classical baselines and larger grids.} Against the classical methods without corrector (bottom block of \Cref{tab:ca_time_poisson}), the router with its best stationary pairing reaches truncation-level accuracy several times faster than multigrid alone (three V-cycles) and than multigrid-preconditioned CG at $128^2$ and $256^2$, and an order of magnitude faster than unpreconditioned or SymGS-preconditioned CG. The reason is the corrector's accuracy on the smooth band: one call removes the modes $|k|\le 31$ to $10^{-6}$ at the price of a single $4096\times4096$ matrix--vector product, whereas a V(2,2)-cycle contracts the whole spectrum by $\approx 0.03$ at the price of four Gauss--Seidel sweeps on the fine grid plus the coarse-grid work. Routing over $\{\mathrm{NO}, \text{multigrid}\}$ is itself about twice as fast as multigrid alone on every grid: the oracle and the router call the corrector once and let the V-cycle remove the rest, which is the case in which the per-unit-cost comparison of \Cref{sec:greedy} matters (the cycle is dearer than the unit). On $256^2$ the router remains the fastest iterative method in every pairing (\caSpHintsBMin--\caSpHintsBMax\ over HINTS, \caSpBestBMin--\caSpBestBMax\ over the best period across the three equations). At $512^2$ the margins shrink but the ordering holds for the stationary pairings, and the $\{\mathrm{NO}, \text{multigrid}\}$ router is the fastest iterative method we measured at this size. With Jacobi and Gauss--Seidel the gains at $512^2$ are small because the corrector's fixed band $|k|\le 31$ leaves the smoother all modes $k \ge 32$, whose slowest member contracts by only $0.98$ per Jacobi sweep at $N=512$ (against $0.71$ at $N=128$): every policy then needs the same $\approx 80$ sweeps after a single corrector call, the wall-clock time is dominated by sweeps, and the schedule barely matters (with Jacobi the router misses $h^2$ within the $8{,}000$-iteration cap on $2$ of $16$ instances, marked in the table). This is the expected limit of a corrector whose band does not grow with the grid: the hybrid's advantage comes from removing the modes the smoother cannot damp, and a corrector with a band that scales with $N$ (at the price of a dense matrix--vector product that grows like the fourth power of the band) or a multigrid member restores the gains, as the multigrid pairing shows.

\begin{table}[H]
\color{red}
\centering
\caption{Poisson: median wall-clock time to relative error $\varepsilon \in \{10^{-3}, h^2, 10^{-8}\}$ on the $128^2$, $256^2$ and $512^2$ grids ($64$, $32$ and $16$ test instances) for every pairing $\{\mathrm{NO}, \text{solver}\}$ and policy, and for the classical baselines without corrector. ``HINTS (best $\tau$)'' is the best of $\tau\in\{5,10,25,50\}$ chosen per cell on the test set itself (the chosen $\tau$ is the subscript). The greedy oracle is the cost-agnostic \Cref{alg:greedy}, the cost-aware oracle is \Cref{alg:greedy} on cost-equalised macro-actions; both use the true error and are not deployable. Bold marks the router where it is the fastest deployable method of its pairing; $\dagger n$: $n$ instances did not reach the tolerance within the iteration cap ($>$: lower bound).}
\resizebox{\textwidth}{!}{\catimepoisson}
\label{tab:ca_time_poisson}
\end{table}

\begin{table}[H]
\color{red}
\centering
\caption{Poisson: paired median speedup of the learned router at $\varepsilon = h^2$ and $10^{-8}$ over HINTS ($\tau{=}25$), over the best fixed period for that cell, over multigrid alone and over multigrid-preconditioned CG, with the one-sided Wilcoxon $p$-value that the router is faster (or, where its median is larger, that it is slower). Bold: $\ge 1.1\times$ with $p < 0.01$.}
\resizebox{\textwidth}{!}{\caspeedpoisson}
\label{tab:ca_speed_poisson}
\end{table}

\begin{figure}[H]
\color{red}
    \centering
    \newfig[0.98\linewidth]{neurips_images/ca_usage_poisson.png}
    \caption{Poisson: fraction of test runs that execute a corrector call at each iteration, for the learned router, the cost-aware oracle and HINTS ($\tau{=}25$), for every pairing (columns) and grid (rows; all test instances of each grid).}
    \label{fig:ca_usage_poisson}
\end{figure}

\begin{table}[H]
\color{red}
\centering
\caption{Poisson: median iterations (and corrector calls) to $\varepsilon = h^2$ for HINTS ($\tau{=}25$), the cost-aware oracle and the learned router on every grid.}
\resizebox{0.9\textwidth}{!}{\causagepoisson}
\label{tab:ca_usage_poisson}
\end{table}

\begin{table}[H]
\color{red}
\centering
\caption{Poisson: measured cost of one classical iteration $c_{\text{solver}}$ and of one corrector call $c_{\mathrm{NO}}$ (residual evaluation included), and the resulting macro-action size $m_{\text{solver}}$ (sweeps per decision; the corrector is always applied once).}
\resizebox{0.9\textwidth}{!}{\cacostspoisson}
\label{tab:ca_costs_poisson}
\end{table}

\begin{table}[H]
\color{red}
\centering
\caption{Poisson, $128\times128$: iteration-based metrics of the main text---relative error after $T = \caT$ iterations and AUC over the $T$ iterations (mean and standard error over $64$ instances; $p$: one-sided paired $t$-test that the method's AUC exceeds the router's). Bold: router better than both solver-only and HINTS.}
\resizebox{0.6\textwidth}{!}{\caaucpoisson}
\label{tab:ca_auc_poisson}
\end{table}

\begin{table}[H]
\color{red}
\centering
\caption{Poisson, $128\times128$: five independent router-training seeds per pairing---mean and standard deviation over seeds of the median work-unit time to $h^2$, the share of test instances on which the seed router's decision sequence to $h^2$ coincides with the main router's, the range over seeds of the paired speedup over HINTS ($\tau{=}25$), and the number of seeds significantly faster ($p<0.01$) than both HINTS ($\tau{=}25$) and the best fixed period.}
\resizebox{0.85\textwidth}{!}{\caseedspoisson}
\label{tab:ca_seeds_poisson}
\end{table}

\begin{figure}[H]
\color{red}
    \centering
    \newfig[0.98\linewidth]{neurips_images/ca_router_decisions_poisson.png}
    \caption{Poisson, $128\times128$: routing decisions (dark: corrector call) of the learned router (top row) and the cost-aware oracle (bottom row) over the first $60$ iterations of every test instance, for every pairing.}
    \label{fig:ca_decisions_poisson}
\end{figure}

\subsection{Convection--diffusion} \label{sec:wc_conv}

The convection--diffusion results (\Cref{tab:ca_time_conv,tab:ca_speed_conv,fig:ca_usage_conv,tab:ca_usage_conv,tab:ca_costs_conv,tab:ca_auc_conv,tab:ca_seeds_conv}) mirror the Poisson ones: at $128^2$ the router is the fastest deployable method in every pairing at $h^2$ and at $10^{-8}$, reproduces the cost-aware oracle, and is $2$--$3\times$ faster than multigrid alone and than multigrid-preconditioned BiCGSTAB; on $256^2$ the same holds in every pairing; at $512^2$ the router is faster than HINTS and than the best fixed period in the Jacobi, GS and SOR pairings and the $\{\mathrm{NO}, \text{multigrid}\}$ router is again the fastest iterative method, with the same mechanism as on Poisson (the Jacobi and GS pairings are limited by the modes outside the corrector's band: the router, the oracle and HINTS all reach $h^2$ after $84$--$85$ Jacobi iterations). The nonsymmetric operator changes only the details: the best fixed period is $\tau{=}10$ rather than $5$ in several cells, the undamped-Jacobi pairing stalls at the near-Nyquist band below $10^{-6}$ for every policy, and GMRES(20) is two orders of magnitude slower than the router.

\begin{table}[H]
\color{red}
\centering
\caption{Convection--diffusion: median wall-clock time to relative error $\varepsilon \in \{10^{-3}, h^2, 10^{-8}\}$ on the $128^2$, $256^2$ and $512^2$ grids for every pairing and policy, and for the classical baselines without corrector (conventions as in \Cref{tab:ca_time_poisson}).}
\resizebox{\textwidth}{!}{\catimeconv}
\label{tab:ca_time_conv}
\end{table}

\begin{table}[H]
\color{red}
\centering
\caption{Convection--diffusion: paired median speedup of the learned router at $\varepsilon = h^2$ and $10^{-8}$ over HINTS ($\tau{=}25$), the best fixed period, multigrid alone and multigrid-preconditioned BiCGSTAB, with one-sided Wilcoxon $p$-values (conventions as in \Cref{tab:ca_speed_poisson}).}
\resizebox{\textwidth}{!}{\caspeedconv}
\label{tab:ca_speed_conv}
\end{table}

\begin{figure}[H]
\color{red}
    \centering
    \newfig[0.98\linewidth]{neurips_images/ca_usage_convdiff.png}
    \caption{Convection--diffusion: corrector usage per iteration of the learned router, the cost-aware oracle and HINTS ($\tau{=}25$) for every pairing (columns) and grid (rows).}
    \label{fig:ca_usage_conv}
\end{figure}

\begin{table}[H]
\color{red}
\centering
\caption{Convection--diffusion: median iterations (and corrector calls) to $\varepsilon = h^2$ on every grid.}
\resizebox{0.9\textwidth}{!}{\causageconv}
\label{tab:ca_usage_conv}
\end{table}

\begin{table}[H]
\color{red}
\centering
\caption{Convection--diffusion: measured per-iteration costs and macro-action sizes.}
\resizebox{0.9\textwidth}{!}{\cacostsconv}
\label{tab:ca_costs_conv}
\end{table}

\begin{table}[H]
\color{red}
\centering
\caption{Convection--diffusion, $128\times128$: iteration-based metrics (conventions as in \Cref{tab:ca_auc_poisson}).}
\resizebox{0.6\textwidth}{!}{\caaucconv}
\label{tab:ca_auc_conv}
\end{table}

\begin{table}[H]
\color{red}
\centering
\caption{Convection--diffusion, $128\times128$: five-seed router-training trials (conventions as in \Cref{tab:ca_seeds_poisson}).}
\resizebox{0.85\textwidth}{!}{\caseedsconv}
\label{tab:ca_seeds_conv}
\end{table}

\begin{figure}[H]
\color{red}
    \centering
    \newfig[0.98\linewidth]{neurips_images/ca_router_decisions_convdiff.png}
    \caption{Convection--diffusion, $128\times128$: routing decisions (dark: corrector call) of the learned router (top row) and the cost-aware oracle (bottom row) over the first $60$ iterations of every test instance.}
    \label{fig:ca_decisions_conv}
\end{figure}

\subsection{Anisotropic diffusion} \label{sec:wc_aniso}

The anisotropic equation is where the two components of the hybrid are most complementary and where the classical solvers are weakest: the point smoothers cannot damp the $x_1$-oscillatory modes at all (\Cref{tab:ca_time_aniso}: none of the five solvers alone reaches $h^2$ within the iteration cap on either grid), whereas the corrector, whose sensor grid keeps full resolution along $x_1$, removes exactly those modes together with the smooth ones in a single call and leaves the smoother only the $x_2$-oscillatory band, which it damps at the isotropic rate. Consequently the learned router reaches $h^2$ in under $1$\,ms at $128^2$ for every pairing---the same absolute time as on the isotropic equations---while HINTS ($\tau{=}25$) needs $3$--$10$\,ms and the best fixed period $1.2$--$2.4$\,ms; at $256^2$ (Jacobi, GS and SOR pairings, $16$ instances) the router needs $9$--$10$\,ms against $15$--$30$\,ms for HINTS and $10$--$13$\,ms for the best period (\Cref{tab:ca_speed_aniso}). The router again matches the cost-aware oracle, its decisions are identical across the five training seeds (\Cref{tab:ca_seeds_aniso}), and it spends the iterations after the single corrector call on the $x_2$-oscillatory band (\Cref{fig:ca_usage_aniso,tab:ca_usage_aniso}). The multigrid and Krylov baselines are not run for this equation (\Cref{sec:wallclock_protocol}).

\begin{table}[H]
\color{red}
\centering
\caption{Anisotropic diffusion: median wall-clock time to relative error $\varepsilon \in \{10^{-3}, h^2, 10^{-8}\}$ on the $128^2$ ($64$ instances, five pairings) and $256^2$ ($16$ instances; Jacobi, GS and SOR pairings) grids for every policy (conventions as in \Cref{tab:ca_time_poisson}; the solver-only runs are capped at $60{,}000$ and $40{,}000$ iterations and never reach $h^2$).}
\resizebox{0.9\textwidth}{!}{\catimeaniso}
\label{tab:ca_time_aniso}
\end{table}

\begin{table}[H]
\color{red}
\centering
\caption{Anisotropic diffusion: paired median speedup of the learned router at $\varepsilon = h^2$ and $10^{-8}$ over HINTS ($\tau{=}25$) and the best fixed period, with one-sided Wilcoxon $p$-values (no classical baselines without corrector are run for this equation).}
\resizebox{0.9\textwidth}{!}{\caspeedaniso}
\label{tab:ca_speed_aniso}
\end{table}

\begin{figure}[H]
\color{red}
    \centering
    \newfig[0.85\linewidth]{neurips_images/ca_usage_anisodiff.png}
    \caption{Anisotropic diffusion: corrector usage per iteration of the learned router, the cost-aware oracle and HINTS ($\tau{=}25$) for every pairing (columns) and grid (rows).}
    \label{fig:ca_usage_aniso}
\end{figure}

\begin{table}[H]
\color{red}
\centering
\caption{Anisotropic diffusion: median iterations (and corrector calls) to $\varepsilon = h^2$.}
\resizebox{0.75\textwidth}{!}{\causageaniso}
\label{tab:ca_usage_aniso}
\end{table}

\begin{table}[H]
\color{red}
\centering
\caption{Anisotropic diffusion: measured per-iteration costs and macro-action sizes.}
\resizebox{0.75\textwidth}{!}{\cacostsaniso}
\label{tab:ca_costs_aniso}
\end{table}

\begin{table}[H]
\color{red}
\centering
\caption{Anisotropic diffusion, $128\times128$: iteration-based metrics (conventions as in \Cref{tab:ca_auc_poisson}).}
\resizebox{0.6\textwidth}{!}{\caaucaniso}
\label{tab:ca_auc_aniso}
\end{table}

\begin{table}[H]
\color{red}
\centering
\caption{Anisotropic diffusion, $128\times128$: five-seed router-training trials (conventions as in \Cref{tab:ca_seeds_poisson}).}
\resizebox{0.85\textwidth}{!}{\caseedsaniso}
\label{tab:ca_seeds_aniso}
\end{table}

\begin{figure}[H]
\color{red}
    \centering
    \newfig[0.98\linewidth]{neurips_images/ca_router_decisions_anisodiff.png}
    \caption{Anisotropic diffusion, $128\times128$: routing decisions (dark: corrector call) of the learned router (top row) and the cost-aware oracle (bottom row) over the first $60$ iterations of every test instance.}
    \label{fig:ca_decisions_aniso}
\end{figure}

\subsection{Solver ensembles} \label{sec:wallclock_ens}

\paragraph{When can a second classical solver help?} The cost-aware oracle is a per-decision comparison of macro-actions, so a router over $\mathrm{NO}\cup\mathcal{W}$ can only beat the best pairwise router of its members if, along the way to the tolerance, the error passes through regimes in which \emph{different} members give the largest reduction per unit cost. We screened this ceiling exhaustively before training any router: for every subset $\mathcal{W}$ of up to three members of $\{$Jacobi, Jacobi (0.67), GS, SymGS, SOR (1.5), multigrid$\}$ (multigrid on the isotropic equations only), every equation and grid, we ran the cost-aware oracle over $\mathrm{NO}\cup\mathcal{W}$ on the test instances and compared its work-unit time to that of the oracle over each member alone (\Cref{tab:ca_ens_screen}). Two facts emerge. At $\varepsilon = h^2$ no ensemble beats the best single solver anywhere: one corrector call followed by one or two macro-actions of the cheapest smoother already reaches truncation-level accuracy, so there is no second regime to exploit. At $\varepsilon = 10^{-8}$, after the corrector has removed the smooth band, the smoother has to damp the whole band $|k| > 31$, whose sub-bands favour different solvers: undamped Jacobi is the cheapest per unit cost on the modes near the band edge (it contracts them by $\cos(2\pi\cdot 32/N)$-like factors that damped Jacobi cannot match), but it does not damp the near-Nyquist modes at all, which damped Jacobi and GS remove in a few sweeps. The oracle over $\{\text{Jacobi}, \text{Jacobi (0.67)}\}$---a \emph{nonstationary} Jacobi iteration whose damping parameter is chosen per macro-action---is therefore faster than every single solver, by $13$--$28\%$ on four of the six (equation, grid) settings at $128^2$ and $256^2$, and adding a third member (GS) changes nothing, i.e., the gain saturates. At $512^2$ multigrid dominates every subset (its cycle removes all sub-bands at once), so the oracle over any ensemble containing it equals the oracle over $\{\mathrm{NO}, \text{multigrid}\}$: adding members neither helps nor hurts. The ``difficulty'' that matters for ensembles is thus not the equation but the width of the spectrum left to the classical members after the corrector, which grows with the grid (fixed band) and with the tolerance.

\paragraph{Nested ensembles.} \Cref{tab:ca_ens_nest} evaluates the learned router over the nested sets $\{\text{Jacobi}\}$, $\{\text{Jacobi (0.67)}\}$, $\{\text{Jacobi}, \text{Jacobi (0.67)}\}$ and $\{\text{Jacobi}, \text{Jacobi (0.67)}, \text{GS}\}$ on every equation at $128^2$ and $256^2$; the single-member routers are the pairwise routers of \Cref{sec:wc_poisson,sec:wc_conv,sec:wc_aniso}, evaluated in the same session on the same instances, and all members keep the macro-action sizes of their pairwise runs. The table reports live and work-unit times to $h^2$, $10^{-6}$ and $10^{-8}$ for the router and its oracle, and the paired speedup at $10^{-8}$ over the best single-member router with its $p$-value (work units are the primary measure here: the nested runs share the machine with other jobs, so their live times are uniformly inflated relative to \Cref{sec:wc_poisson,sec:wc_conv,sec:wc_aniso}, while work units are determined by the decision sequence alone). Two properties are visible. \emph{Monotonicity}: the router over a larger set is never slower than the best of its members beyond timer noise, because the oracle ignores a member that never gives the largest reduction per unit cost and the router imitates it (adding GS to $\{\text{Jacobi}, \text{Jacobi (0.67)}\}$ leaves the decision sequence unchanged); the earlier ensembles of \Cref{tab:ca_ens} show the same for four- and five-member sets at $h^2$. \emph{Strict separation}: at $10^{-8}$ the two-member router is significantly faster than every single-member router in \caNestWins\ of the \caNestNum\ ensemble cells (paired Wilcoxon $p<0.01$; speedups \caNestSpMin--\caNestSpMax\ over the best single solver), reproducing its oracle, and it is never significantly slower. The learned policy is the nonstationary damping schedule described above: undamped Jacobi macro-actions while the band-edge modes dominate the residual, damped Jacobi once the near-Nyquist modes do, which the router detects from the decay rate of the residual over the last macro-action.

\begin{table}[H]
\color{red}
\centering
\caption{Oracle-level screening of ensembles: for every equation and grid, the fastest single solver and the fastest ensemble of up to three members under the cost-aware oracle (median work-unit time over $16$ test instances), at $\varepsilon = h^2$ and $10^{-8}$; the ratio is the speedup of the ensemble over the best single solver, and the last column is the ratio for $\{\text{Jacobi}, \text{Jacobi (0.67)}\}$.}
\resizebox{\textwidth}{!}{\caensscreen}
\label{tab:ca_ens_screen}
\end{table}

\begin{table}[H]
\color{red}
\centering
\caption{Nested ensembles: median time to $\varepsilon \in \{h^2, 10^{-6}, 10^{-8}\}$ of the learned router (live and work units) and of the cost-aware oracle (work units) over $\mathrm{NO}\cup\mathcal{W}$, for the single members and the nested two- and three-member sets, all evaluated in one session on the same instances ($64$ at $128^2$, $32$ at $256^2$, $16$ for anisotropic diffusion at $256^2$). The last columns give the paired speedup at $10^{-8}$ over the best single-member router (resp.\ oracle) with the one-sided Wilcoxon $p$-value; bold: $\ge 1.1\times$ with $p<0.01$.}
\resizebox{\textwidth}{!}{\caensnest}
\label{tab:ca_ens_nest}
\end{table}

\paragraph{Larger ensembles at truncation level.} \Cref{tab:ca_ens} repeats the ensembles of \Cref{sec:experiments} ($\{\text{Jacobi}, \text{GS}\}$ up to all five solvers) at $128^2$ and $\varepsilon = h^2$, where the screening predicts no gain: the ensemble router matches the best pairwise router of its members to within about $\pm 20\%$ in live time (ratio of medians \caEnsVsPairMin--\caEnsVsPairMax\ over the \caNumEns\ cells) and within $1.5\%$ in work units, without knowing in advance which member is fastest, and is \caEnsVsSolverMin--\caEnsVsSolverMax\ faster than the fastest member alone; \Cref{tab:ca_ensusage} lists how its iterations are distributed over the members. Its oracle is up to $25\%$ faster in a few cells because, at $h^2$, several cost-equalised smoothers are near-ties and the oracle's choice among them depends on the spectral content of the current error, which the router's residual-norm features reflect only indirectly (agreement with the oracle $55$--$83\%$ on its own rollouts against $88$--$99\%$ for pairwise routers). More capacity or imitation data does not change this: routers with twice the hidden width, three times the DAgger rounds and twice the oracle rollouts (\Cref{tab:ca_ens_big}) take the same number of iterations and corrector calls to $h^2$ as the default routers on \caEnsBigSameMin--\caEnsBigSameMax\ of the test instances and are significantly faster in \caEnsBigNFaster\ and significantly slower in \caEnsBigNSlower\ of the \caNumEnsBig\ cells ($p<0.05$); the residual disagreements are near-ties whose resolution does not affect the time to $h^2$. Deciding \emph{when} to call the corrector---the decision that dominates the wall-clock gains---is learned equally well in the ensembles.

\begin{table}[H]
\color{red}
\centering
\caption{Ensembles of \Cref{sec:experiments} at $128\times128$ and $\varepsilon = h^2$: median live time of the best member solver alone, of the best pairwise router among the members (paired speedup of the ensemble router over it in parentheses), and of the router and the cost-aware oracle over $\mathrm{NO}\cup\mathcal{W}$ ($64$ instances; pairwise baselines from the pairwise runs on the same instances).}
\resizebox{\textwidth}{!}{\caens}
\label{tab:ca_ens}
\end{table}

\begin{table}[H]
\color{red}
\centering
\caption{Ensembles at $128\times128$: mean (standard deviation over instances) fraction of the iterations up to the $h^2$ crossing that the ensemble router spends in each member and in the corrector.}
\resizebox{0.9\textwidth}{!}{\caensusage}
\label{tab:ca_ensusage}
\end{table}

\begin{table}[H]
\color{red}
\centering
\caption{Ensembles at $128\times128$: one-sided Wilcoxon $p$-values at $\varepsilon = h^2$ for the ensemble router against the fastest member solver alone (censored runs at their lower bound), against the best pairwise router in both directions, and against the cost-aware oracle over the same ensemble (alternative: the ensemble router is slower).}
\resizebox{0.9\textwidth}{!}{\castatsens}
\label{tab:ca_stats_ens}
\end{table}

\begin{table}[H]
\color{red}
\centering
\caption{Control on router capacity for the ensembles ($128\times128$, $64$ instances): the default ensemble router (hidden width $64$, $2$ DAgger rounds, $128$ oracle rollouts) against one with hidden width $128$, $6$ DAgger rounds and $256$ rollouts. Times to $h^2$ are median work units net of the per-decision cost; ``same decisions'' is the share of instances on which both routers take the same number of iterations and corrector calls to $h^2$; $p$ is the one-sided Wilcoxon $p$-value on the paired net times; ``agreement'' is each router's agreement with the oracle on its own rollouts during training.}
\resizebox{\textwidth}{!}{\caensbig}
\label{tab:ca_ens_big}
\end{table}

\subsection{Overheads and amortisation} \label{sec:wallclock_overheads}

\Cref{tab:ca_overheads} lists the measured cost of every operation as a function of $N$, including one decision of the LSTM router of \Cref{sec:experiments} (full-state input of size $4N^2$, hidden size 256, 4 layers): at $128^2$ one LSTM decision costs \caLstmMs\,ms, about \caLstmOverJacobi\ Jacobi sweeps, which is why the wall-clock study uses the scalar-feature router ($\approx 5\,\mu$s, independent of $N$). The learned router's decisions therefore account for well under $1\%$ of its solve time, and its corrector calls for $50$--$70\%$ of the time to $h^2$ (one or two calls). \Cref{tab:ca_amort} accounts for training: the corrector (shared with HINTS) is fit in minutes, the router in seconds to a few minutes per pairing, and the break-even number of solves after which the router's own training cost is recovered relative to HINTS ($\tau{=}25$) with the same corrector is in the thousands at $128^2$ and in the hundreds on the larger grids.

\begin{table}[H]
\color{red}
\centering
\caption{Measured single-thread cost of one operation as a function of the grid size (Poisson; median of repeated calls). Multigrid uses a direct solve below $64^2$ and is not listed there.}
\resizebox{0.9\textwidth}{!}{\caoverheads}
\label{tab:ca_overheads}
\end{table}

\begin{table}[H]
\color{red}
\centering
\caption{Training-cost amortisation for the GS pairing: corrector training (data generation and least-squares fit, CPU), router training (oracle rollouts and DAgger rounds), median wall-clock saved per solve at $\varepsilon = h^2$ relative to HINTS ($\tau{=}25$) and to the solver alone, and the number of solves after which the router's training cost is recovered relative to HINTS.}
\resizebox{\textwidth}{!}{\caamort}
\label{tab:ca_amort}
\end{table}

\subsection{Verification of the theoretical assumptions} \label{sec:wallclock_assumptions}

\Cref{th:suboptgreedy} and Propositions~\ref{th:weaklyalphasupermodular}--\ref{th:simultaneoussupermodular} assume that every error propagation map $G_j = I - C_j\circ\mathcal{L}_h^a$ is Lipschitz with constant $\rho_j < 1$ (respectively $\le 1$ with a shared eigenbasis), zero-preserving, and---for postfix monotonicity---invertible; \Cref{th:consistency} assumes that the per-step errors $\tilde c_j$ are bounded above and that at least one solver cannot annihilate the error in one step. \Cref{tab:ca_assump} verifies these conditions numerically for every operation of the study on the mean-free subspace (the constant mode is the null space of the periodic operator and is excluded from the data and from every iteration): the Euclidean Lipschitz constant $\|G_j\|_2$ (power iteration on $G_j^{\top}G_j$; exact from the Fourier symbol for Jacobi), the Lipschitz constant $\|G_j\|_A = \|A_s^{1/2} G_j A_s^{-1/2}\|_2$ in the energy norm of the symmetric part $A_s$ of the discrete operator, the same for the macro-action $G_j^{m_j}$, the spectral radius, the smallest observed contraction $\sigma_{\min} = \min\|G_j x\|/\|x\|$, $\|C_j(0)\|$, and the commutator $\|G_j G_{\mathrm{NO}} - G_{\mathrm{NO}} G_j\|$ on random vectors. \Cref{tab:ca_assump_paths} evaluates the quantities of \Cref{th:suboptgreedy} and \Cref{th:consistency} along the cost-aware oracle's own path to $\varepsilon = 10^{-8}$ (a horizon of tens of macro-actions; to $h^2$ the path has one or two steps).

\paragraph{What holds, and in which norm.} (i) Every classical solver is zero-preserving exactly and contracts in the sense that matters for convergence: the spectral radius of $G_j$ is below $1$ in every setting (at most \caRhoSpecMax), and so is the energy-norm Lipschitz constant of the solver and of its macro-action (at most \caRhoAMax). The \emph{Euclidean} Lipschitz constant, however, exceeds $1$ for Gauss--Seidel, SOR and symmetric GS (up to \caRhoTwoGsMax\ for SOR), whose iteration matrices are non-normal; the proofs of Propositions~\ref{th:weaklyalphasupermodular}--\ref{th:simultaneoussupermodular} use only the triangle inequality and the composition rule for Lipschitz constants, so they hold verbatim with $\|\cdot\|_2$ replaced by any norm, and in the energy norm the Lipschitz condition is satisfied by all classical solvers in all settings. For (damped) Jacobi both constants agree (the maps are Fourier-diagonal); the values are close to $1$ because the slowest modes (the smoothest ones, and the near-Nyquist mode for undamped Jacobi) are barely damped by a point smoother---which is exactly why the corrector is needed. (ii) The corrector's map is zero-preserving, contracts every band mode to $\caBandMax$ or better (on average to $10^{-4}$--$10^{-6}$) and leaves the modes outside the band untouched, so its Lipschitz constant is exactly $1$ (not $<1$): \Cref{th:weaklyalphasupermodular} in its strict form does not apply to the corrector, but \Cref{th:simultaneoussupermodular} does, and its second condition---a shared eigenbasis---holds exactly for the ensembles $\{\mathrm{NO}, \text{Jacobi}\}$, $\{\mathrm{NO}, \text{Jacobi (0.67)}\}$ and $\{\mathrm{NO}, \text{Jacobi}, \text{Jacobi (0.67)}\}$ (all maps are diagonal in the Fourier basis; measured commutators $\approx 10^{-6}$, at the corrector's float32 precision), so $h$ is supermodular for them and \Cref{th:suboptgreedy} holds with $\alpha = 1$; for GS, SOR, symmetric GS and multigrid the commutators are small but nonzero ($10^{-3}$--$10^{-1}$). (iii) Invertibility fails for most maps in exact arithmetic: a Gauss--Seidel, SOR or symmetric-GS sweep annihilates an error concentrated at the first grid point ($\sigma_{\min} = 0$), undamped Jacobi annihilates the mode $(N/4, N/4)$, and the corrector annihilates the band; postfix monotonicity is therefore not guaranteed, which \Cref{sec:greedy} anticipates (``our experiments show that greedy routers remain effective even when invertibility is not met''). (iv) The worst-case ratio $\alpha(O)$ of \Cref{th:weaklyalphasupermodular}, which only uses the Lipschitz constants, is loose here: along the oracle's path to $10^{-8}$ (\Cref{tab:ca_assump_paths}) $\sum_i \rho_{O_i}^2$ is close to $T$ because each $\rho$ is dominated by the slowest mode, so the bound ranges from $\caAlphaBoundMin$ (the multigrid pairing, whose cycle contracts everything) to several hundred, and is infinite on the few instances where $\sum_i \rho_{O_i}^2 \ge T$. The empirically measured supermodularity ratio of \Cref{eqn:alpha_supermod} along the same path---with $S$ the oracle's prefix and $S'$ its suffix---is at most \caAlphaHatMax\ (median \caAlphaHatMed) in every setting, and at most \caAlphaHatMaxH\ on the short paths to $h^2$: the objective behaves supermodularly on the states the greedy rule actually visits, which is the regime in which \Cref{th:suboptgreedy} gives the $(1-e^{-1})$ guarantee. (v) For \Cref{th:consistency}, the per-step errors are bounded by the current error ($\max_j \tilde c_j / \|e^{(t)}\|^2 \le 1$ at every visited state) and at least one solver leaves a positive fraction of it ($\min_t \max_j \tilde c_j / \|e^{(t)}\|^2 > 0$), so the surrogate's Bayes consistency applies.

\begin{table}[H]
\color{red}
\centering
\caption{Numerical verification of the assumptions of Propositions~\ref{th:weaklyalphasupermodular}--\ref{th:simultaneoussupermodular} for every operation of the study (mean-free subspace). $\|\cdot\|_2$: Euclidean Lipschitz constant of the error propagation map; $\|\cdot\|_A$: Lipschitz constant in the energy norm of the symmetric part of $\mathcal{L}_h$, for one iteration and for the macro-action of $m_j$ iterations; $\rho$: spectral radius; $\sigma_{\min}$: smallest observed contraction (a unit error at the first grid point for the triangular sweeps, the minimum over Fourier modes otherwise; for the corrector, the largest remaining fraction of a band mode); $\|C_j(0)\|$: zero preservation; last column: commutator with the corrector's map on random unit vectors (Proposition~\ref{th:simultaneoussupermodular}). For (damped) Jacobi the constants are exact from the Fourier symbol (the value with the Nyquist modes included; without them undamped Jacobi has $\|G\|_2 = \rho = 0.9994$ at $128^2$).}
\resizebox{\textwidth}{!}{\caassump}
\label{tab:ca_assump}
\end{table}

\begin{table}[H]
\color{red}
\centering
\caption{Quantities of \Cref{th:suboptgreedy}, Proposition~\ref{th:weaklyalphasupermodular} and \Cref{th:consistency} along the cost-aware oracle's path to $\varepsilon = 10^{-8}$ ($8$ test instances; medians unless stated): number of macro-actions $T$, $\sum_i\rho_{O_i}^2$ with the energy-norm constants of the chosen macro-actions, the resulting $\alpha(O)$ of Proposition~\ref{th:weaklyalphasupermodular} (median over the instances on which $T - \sum_i\rho_{O_i}^2 > 0$; $\infty$ otherwise), the empirical supermodularity ratio $\hat\alpha$ of \Cref{eqn:alpha_supermod} evaluated with $S$ the oracle's prefix and $S'$ its suffix (maximum and median over prefixes and instances), and the normalised bounds of \Cref{th:consistency}: the largest ratio $\max_j \tilde c_j/\|e^{(t)}\|^2$ over visited states ($\bar E$) and the smallest value of $\max_j \tilde c_j/\|e^{(t)}\|^2$ ($E_{\min}$).}
\resizebox{\textwidth}{!}{\caassumpB}
\label{tab:ca_assump_paths}
\end{table}
